\section{Choose when pixels become public output}
For the completed memory interface, validate the entire stream before changing the
caller's output. Validation reconstructs the previous pixel and cache privately,
counts the pixels, checks every command extent, and confirms the final marker. It
publishes a small metadata descriptor only after all those checks succeed.

The decoder calls this validation stage into private metadata, then checks framebuffer
capacity. Only then does it reset the reconstruction state and make a second pass that
stores pixels. A malformed stream therefore leaves the framebuffer unchanged, and a
malformed stream with insufficient capacity reports the format error first. This is
the same publication principle used in the PPM chapter, applied to a stateful stream.

During the accepted second pass, each word store can update the connected display.
That observable construction order does not mean the decoder accepted an unvalidated
prefix. It already checked the complete file before the first public pixel store.
An IDE may combine repaint requests, so a screenshot need not show a separate rendered
frame for every store. Debugger state and memory observations identify the individual
instruction boundaries.

Two passes perform more reconstruction work than one. Record that cost when comparing
implementations. A future streaming interface could allow a valid prefix on failure,
but that would be a different published failure promise. An optimization of this
interface must retain its unchanged-output guarantee.
